\section{Solutions Explorées et Abandonnées}

Cette section documente brièvement les pistes techniques explorées mais finalement écartées. Elle sert de guide aux futurs développeurs pour éviter de répéter les mêmes essais infructueux.

\subsection{Régulateur de Tension Centralisé (Traco Power)}
Le plan initial prévoyait un unique module Traco Power pour alimenter tout le système à partir d'une seule batterie LiPo. En pratique, les pics de courant des moteurs lors des accélérations provoquaient des chutes de tension sur le rail partagé, causant des redémarrages intempestifs de la Raspberry Pi (brown-outs). La solution retenue a été d'isoler électriquement les domaines d'alimentation : la Raspberry Pi est alimentée par une powerbank USB-C indépendante, le système de traction par une batterie LiPo principale, et le mécanisme de montée/descente de la sonde par une seconde batterie LiPo moins puissante, dédiée uniquement à cet actionneur.

\subsection{Vision YOLO pour la Navigation}
L'exécution de modèles YOLO sur la Raspberry Pi 4 n'atteignait que 2--3 images par seconde, insuffisant pour un contrôle réactif. La vision haute résolution a donc été remplacée par une segmentation colorimétrique HSV à 160x120 pixels avec saut de trames, atteignant 10--15 images par seconde. YOLO reste pertinent pour la détection de végétaux et de pestes, mais nécessite un matériel plus puissant (voir section~\ref{sec:ameliorations}).

\subsection{micro-ROS sur Arduino}
La bibliothèque micro-ROS aurait permis d'intégrer l'Arduino directement dans le graphe de communication ROS 2. Cependant, les ressources limitées de l'Arduino Uno (2~Ko de RAM, 32~Ko de Flash) causaient des instabilités mémoire lors de l'exécution simultanée de la sécurité ultrasonique et du client micro-ROS. L'équipe a préféré maintenir un protocole propriétaire léger sur l'Arduino, la logique ROS 2 restant entièrement côté Raspberry Pi.

% \subsection{IMU BNO055}
% L'IMU BNO055 avait été envisagée pour corriger les dérives de trajectoire dues au glissement (skid-steering). En pratique, le champ magnétique des moteurs perturbait la boussole interne, rendant les données de cap peu fiables. L'odométrie visuelle par centroïde de ligne, couplée au LiDAR, a suffi pour les missions de suivi de rang.

\subsection{Navigation Globale avec Nav2}
La stack Nav2 de ROS 2, qui inclut des planificateurs de trajectoire de type A*, a été étudiée. Son intégration exige une carte de l'environnement (via SLAM) et un système de localisation robuste, deux éléments non disponibles dans le périmètre du projet. Le contrôleur PD réactif, bien que plus simple, a été suffisant pour les missions en rang droit.
